\documentclass{ieeeaccess}
\usepackage{cite}
\usepackage{amsmath,amssymb,amsfonts}
\usepackage{algorithm}
\usepackage{algorithmic}
\usepackage{graphicx}
\usepackage{float}
\usepackage{textcomp}
\usepackage{hyperref}
\usepackage{enumitem}
\usepackage{caption}

\def\BibTeX{{\rm B\kern-.05em{\sc i\kern-.025em b}\kern-.08em
    T\kern-.1667em\lower.7ex\hbox{E}\kern-.125emX}}
\renewcommand{\thetable}{\Roman{table}}
\begin{document}

\title{A Semantic-Aware Spatio-Temporal Data Warehouse for Micro-Mobility:
The Case of the MakFleet Bodaboda Network at Makerere University}

\author{\uppercase{Ddumba Jonah}\authorrefmark{1},
\uppercase{Ssebwana John Paul}\authorrefmark{1},
\uppercase{Sebunya Ronaldo}\authorrefmark{1},
and \uppercase{Namuli Evelyn}\authorrefmark{1}}
\address[1]{Department of Computer Science, Makerere University,
Kampala, Uganda}

\tfootnote{This work was developed as part of the Data Warehousing
and Business Intelligence coursework submitted in April 2026.}


\begin{abstract}
The emergence of Internet of Things telematics in micro mobility presents notable opportunities for optimizing logistics on college campuses. However, conventional relational database technologies such as those based on the Star or Snowflake schema design, lack the ability to cope with high-velocity, multimodal sensor streams and erratic spatio-temporal dynamics of motorcycle taxis also known as bodabodas. This paper proposes a novel semantic aware data warehousing architecture tailored for the MakFleet motorcycle taxiing service at Makerere University. This envisions a Spatio-Temporal Knowledge Graph which integrates high frequency telemetry data including global positioning system and accelerometer measurements with the topological constraints and safety regulations of the Makerere campus in mind. Through the adoption of semantic validation this model achieves epistemic trust in the ingested data guaranteeing its authenticity and provenance. The structured semantics form a critical data foundation for machine learning algorithms in artificial intelligence applications allowing to detect real-time anomalies associated with reckless driving and perform more accurate predictions for allocating resources in the fleet.
\end{abstract}

\begin{keywords}
Bodaboda Telemetry, Epistemic Trust, Graph Representation Learning,
MakFleet, Micro-Mobility, Semantic Web, Spatio-Temporal Data Warehouse.
\end{keywords}

\titlepgskip=-15pt
\maketitle
\sloppy


\section{INTRODUCTION}

The micro mobility landscape at Makerere University relies heavily on the motorcycle taxi network for rapid intra-campus transit. Through the
MakFleet initiative, these motorcycles are increasingly equipped with Internet of Things telematics devices that stream high-frequency
data regarding location, acceleration and engine idling
\cite{gubbi2013,chen2017}. However, traditional Online Transaction Processing  systems and legacy relational data warehouses are fundamentally ill-equipped to handle the chaotic and informal routing
behaviors characteristic of motorcycle taxis also known as bodabodas in Uganda \cite{kimball2013}. Due to the fact that  bodabodas do not strictly adhere to mapped roads, they often navigate pedestrian pathways and complex campus topologies
erratically.The raw telemetry data they produce remains analytically insufficient without
a structured semantic framework to contextualize it against spatial
geography and safety policies.

This paper proposes a paradigm shift through the construction of a Spatio-Temporal Data Warehouse that utilises semantic web technologies
to transform raw MakFleet Internet of Things streams into scientifically valid evidence. By moving far  beyond mere fleet tracking, this research provides a foundational knowledge infrastructure capable of supporting advanced causal inference and predictive artificial intelligence models for campus safety.


\section{LITERATURE REVIEW}
\label{sec:literature}


The dominant methodology of data warehousing has long been anchored in dimensional relational models, most notably the Kimball Star Schema.
While these models excel at structured, batch oriented reporting, they are structurally ill-suited to the fluid dynamics of mobility data generated by automobiles \cite{kimball2013,chen2012}. Business Intelligence frameworks that are  designed for tabular big data often fail to maintain analytical coherence when queries require multi-dimensional spatio-temporal joins therefore  creating significant bottlenecks under real-time streaming requirements \cite{liang2023survey}.

Spatio-temporal data warehouses have been proposed to address the need for extensible data types that go beyond the native representations of Structured Query Language  schemas \cite{pelekis2014}. Schema rigidity has been noted as a primary challenge that prevents relational systems from handling high-velocity, multimodal data \cite{hamdi2022}. Whereas  traditional relational databases can support temporal queries, fundamental scalability limits emerge when spatial graph topology must be co-modelled alongside time series particularly in complex environments such as the MakFleet network  adding complexity to the implementation\cite{pelekis2014}.

The increased use of Global Positioning System enabled devices and Internet of Things sensors has produced an exponential growth in mobility data streams in the 21st century \cite{gubbi2013}. The raw sensor output become analytically meaningful only when coupled with semantic enrichment layers \cite{gubbi2013}. Multi-layered semantic
frameworks using knowledge graphs for Internet of Things ecosystems have achieved high
interoperability through edge processing demonstrating that ontology-driven architectures are operationally superior to standard
relational pipelines \cite{pan2024}.

The integration of Large Language Models into the knowledge graph construction pipeline represents the frontier of semantic data engineering.These Large Language Models can automate entity recognition and relation extraction while knowledge graphs ground Large Language Model outputs in structured verifiable facts \cite{pan2024}.

Standard tabular machine learning models are demonstrably inadequate for transport anomaly detection because they cannot encode the
topological structure of road networks. This leaves the Graph-level anomaly detection which
outperforms node level approaches for identifying complex traffic anomalies such as congestion and dangerous driving \cite{zhu2022}.
This finding directly informs the selection of the Spatio-Temporal Graph Neural Network (ST-GNN) architecture in the MakFleet engine. For
temporal modelling, the sinusoidal positional encoding introduced by the Transformer architecture has become the standard mechanism for injecting temporal order into sequence models \cite{vaswani2017}.

Privacy vulnerabilities in mobility data systems specifically the risk that Global Positioning System trajectory datasets can be re-identified despite anonymisation are well established and taken into consideration  \cite{demontjoye2013}. Spatial $k$-anonymity and differential privacy noise injection are the primary mitigations endorsed in the literature reviewed\cite{demontjoye2013}.


\section{COMPUTATIONAL PROBLEM STATEMENT}

To ensure rigorous system design, the MakFleet analytical engine is formalised using the 8 layer locator framework. This is defined as an
\textbf{(L0)} operational safety and spatio-temporal optimisation
problem utilising \textbf{(L1)} weak supervision from Internet of Things telematics
and geospatial mapping to solve an \textbf{(L2)} anomaly detection (reckless driving) and temporal forecasting task on \textbf{(L3)} multivariate time series and spatial graph data, trained under an \textbf{(L4)} continual learning regime that adapts to shifting student traffic patterns between academic semesters, under strict \textbf{(L5)} privacy-preserving and out-of-distribution (OOD) robustness constraints deployed via \textbf{(L6)} edge-to-cloud streaming into a university-governed semantic warehouse utilising
\textbf{(L7)} Spatio-Temporal Graph Neural Networks (ST-GNN) and rule-based symbolic model families.

This formalisation dictates a departure from standard relational
architectures. Because the data geometry (L3) involves highly interconnected spatial routes and sequential time-series, a standard
tabular Star Schema suffers from performance bottlenecks due to complex Structured Query Language joins. Instead, the required inductive bias (L7) necessitates a Knowledge Graph infrastructure that natively supports message passing across spatial nodes ensuring that anomaly detection (L2) accurately flags reckless driving based on contextual topology.


\section{METHODOLOGY: SEMANTIC DATA ENGINEERING}
\label{sec:methodology}

A semantic data engineering pipeline is utilised to handle high-frequency telemetry data. The pipeline operates in three continuous phases.

\textbf{Internet of Things Ingestion and Provenance:} Raw telemetry is ingested via
edge-to-cloud streaming from sensors attached to the bodabodas. Cryptographic hashing is applied at ingestion to establish unbroken data lineage and epistemic trust.

\textbf{Semantic Map Matching:} Raw Global Positioning System coordinates are snapped to the
Makerere University spatial ontology using semantic map-matching algorithms distinguishing between formal road networks and informal
pathways.

\textbf{Event Enrichment:} High-frequency accelerometer data is aggregated into discrete semantically meaningful events and linked to
the specific spatial edge where the event occurred.

Traditional relational architectures including Star and Snowflake
schemas, are fundamentally incapable of modelling the complex dependencies inherent in the L3 data geometry of the MakFleet system.
Because the data involves high velocity multivariate time series mapped
onto a highly interconnected spatial graph specifically the Makerere University campus map representing these relationships in a tabular
format would require expensive and multi-way recursive SQL joins to resolve
spatial proximity hence why not utilised. Under the (L6) systems constraints of real-time
edge-to-cloud streaming, the computational overhead of these joins leads to significant latency rendering the system ineffective for immediate safety interventions. Furthermore, a rigid relational schema cannot accommodate the multimodal missingness or the semantic drift common in informal transport networks. Consequently, the adoption of an (L7)
Ontology-Driven Knowledge Graph is required. Unlike static relational
tables, this architecture natively supports message passing across
spatial nodes allowing the model to incorporate topological context such as the distinction between a paved road and an informal pedestrian
shortcut directly into the inference process to ensure scientifically valid and governable artificial intelligence evidence.


\section{DATA MINING AND MODEL ARCHITECTURES}
\label{sec:model_arch}
To perform advanced anomaly detection and predictive fleet allocation,
the analytical engine leverages a custom Spatio-Temporal Graph Neural
Network (ST-GNN). ST-GNNs have the inductive bias required to process
relational and temporal mobility data simultaneously. In the proposed
architecture, the Makerere campus acts as the graph structure, with
each intersection as a node and each road segment as an edge. The
ST-GNN uses message passing to aggregate features from neighbouring
spatial nodes over time, enabling the model to differentiate between a
contextually safe stop near a known zebra crossing and true reckless
driving such as harsh braking on an open road.

As illustrated in Fig.~\ref{fig:architecture_diag}, the architecture
consists of two Graph Convolutional Network (GCN) encoder layers
followed by a per-node Long Short-Term Memory (LSTM) temporal
aggregator. This neuro-symbolic design---combining the GCN spatial
encoder with rule-based semantic event enrichment---directly satisfies
the (L7) inductive bias requirement and propagates risk signals across
connected campus pathways.

\begin{figure}[H]
\centerline{\includegraphics[width=0.55\columnwidth,keepaspectratio]%
  {architecture_diag.png}}
\caption{ST-GNN architecture interacting with the Spatio-Temporal
Knowledge Graph.}
\label{fig:architecture_diag}
\end{figure}

\subsection{Mathematical Formulations}

The spatial encoding at each graph layer follows the GCN propagation
rule of Kipf and Welling \cite{kipf2017}:
\begin{equation}
\mathbf{H}^{(l+1)} = \sigma\!\left(
  \tilde{\mathbf{D}}^{-\frac{1}{2}}
  \tilde{\mathbf{A}}\,
  \tilde{\mathbf{D}}^{-\frac{1}{2}}
  \mathbf{H}^{(l)}\mathbf{W}^{(l)}
\right)
\end{equation}
where $\tilde{\mathbf{A}} = \mathbf{A} + \mathbf{I}_N$ is the
adjacency matrix with added self-loops, $\tilde{\mathbf{D}}$ is its
degree matrix, $\mathbf{W}^{(l)}$ is a trainable weight matrix, and
$\sigma$ is the ReLU activation. Two such layers are applied per
snapshot, producing node embeddings
$\mathbf{H}_t \in \mathbb{R}^{N \times 64}$.

The temporal aggregator stacks $T=4$ consecutive snapshot embeddings
and processes them through a per-node LSTM:
\begin{equation}
\mathbf{h}_t,\,\mathbf{c}_t =
  \text{LSTM}\!\left(\mathbf{H}_t,\,\mathbf{h}_{t-1},\,
  \mathbf{c}_{t-1}\right), \quad
\hat{p}_i = \sigma_s\!\left(
  \mathbf{W}_2\,\text{ReLU}(\mathbf{W}_1\,\mathbf{h}_{T,i}+\mathbf{b}_1)
  +\mathbf{b}_2\right)
\end{equation}
where $\mathbf{h}_{T,i}$ is the final LSTM hidden state for node $i$
and $\sigma_s$ is the sigmoid function producing an anomaly probability
$\hat{p}_i \in [0,1]$.

To address the class imbalance in the MakFleet dataset, a
Weighted Binary Cross-Entropy (WBCE) loss is applied:
\begin{equation}
\mathcal{L} = -\frac{1}{N}\sum_{i=1}^{N}
  \Bigl[w^{+}\,y_i\log(\hat{p}_i)
  + (1-y_i)\log(1-\hat{p}_i)\Bigr]
\end{equation}
where $y_i \in \{0,1\}$ is the ground-truth anomaly label and
$w^{+} = 19.0$ is the positive-class weight calibrated to the
approximately 10:1 normal-to-anomaly ratio in the dataset. This
formulation forces the model to penalise missed anomalies
proportionally more than false positives, which is the operationally
correct priority for campus safety.

To navigate the semantic warehouse, a Causal Subgraph Retrieval
workflow is implemented that enables administrators to query the
Knowledge Graph for structured evidence whenever an anomaly is flagged.
Algorithm~\ref{alg:retrieval} formalises this process.

\begin{algorithm}[!t]
\caption{Causal Subgraph Retrieval for Anomaly Evidence}
\label{alg:retrieval}
\begin{algorithmic}[1]
\STATE \textbf{Input:} Flagged event identifier $e_{\text{id}}$,
       Knowledge Graph $G$
\STATE \textbf{Output:} Causal subgraph $S$ with provenance chain
\STATE \textit{Anchor Lookup:} Retrieve node
       $e \leftarrow G.\texttt{TelematicsEvent}(e_{\text{id}})$.
\STATE \textit{Location Traversal:} Follow
       $e \xrightarrow{\texttt{LOCATED\_AT}} i$ to obtain the
       flagged intersection $i$.
\STATE \textit{Driver Context:} Follow
       $e \xrightarrow{\texttt{CAUSED\_BY}} d \xrightarrow{\texttt{DRIVES}} v$
       to obtain driver and vehicle nodes.
\STATE \textit{Neighbourhood Expansion:} Collect all intersections
       reachable from $i$ via
       $\texttt{ADJACENT\_TO}^{*1..2}$ (2-hop ego-graph).
\STATE \textit{History Retrieval:} Retrieve prior anomalous events
       $\{e'\}$ linked to $d$ via \texttt{CAUSED\_BY}.
\STATE \textbf{return} Subgraph $S = \{e, i, d, v,
       \text{neighbours}, \{e'\}\}$ with full provenance.
\end{algorithmic}
\end{algorithm}

\subsection{Model Configuration}
The ST-GNN is configured with \textbf{6 node input features}: latitude,
longitude, betweenness centrality, degree centrality, event count per
15-minute window, and mean vehicle speed. The spatial encoder consists
of two \texttt{GCNConv} layers with hidden dimension 64 followed by
ReLU activations. The temporal aggregator is a 2-layer LSTM with hidden
dimension 64 operating on sequences of $T=4$ consecutive 15-minute
graph snapshots. The output head is a two-layer MLP with sigmoid
activation producing a per-node binary anomaly probability. The model
is optimised with Adam ($\text{lr}=10^{-3}$) and a Cosine Annealing
learning rate schedule.


\section{DATASET AND KNOWLEDGE GRAPH DESCRIPTION}
\label{sec:dataset}

The MakFleet dataset is a synthetically simulated telemetry corpus
generated over the actual Makerere University road network retrieved
from OpenStreetMap via OSMnx \cite{boeing2017}. Because no anonymised
public dataset exists for bodaboda operations in Uganda, simulation
over real campus topology is the standard methodology in
privacy-constrained mobility research. The simulation models
\textbf{15 bodaboda vehicles}, each equipped with a virtual IoT
telematics unit, across two academic semesters: Semester~1
(January--April 2024) and Semester~2 (August--October 2024). In total,
\textbf{826,260 telemetry records} are generated at \textbf{5-second
intervals}. Each record contains GPS coordinates ($\pm$5~m Gaussian
noise), speed (km/h), 3-axis accelerometer readings (m/s$^2$), engine
state, and heading. The campus road network comprises \textbf{212
intersections} and \textbf{476 road segments}. High-risk drivers
(risk profile $>0.7$) have anomalous events injected at calibrated
rates (5\% harsh braking, 3\% rapid acceleration), yielding an overall
anomaly rate of 8.8\%.

These records are unified into a Spatio-Temporal Knowledge Graph (ST-KG)
stored in Neo4j. The node and relationship schemas are summarised in
Table~\ref{tab:kg_entities} and Table~\ref{tab:kg_relationships},
respectively. A key ontological commitment of the ST-KG is the
distinction between \texttt{TelematicsEvent} and
\texttt{ContextualStop} as \emph{distinct node labels}: a
\texttt{MATCH (n:TelematicsEvent)} query will never return a
\texttt{ContextualStop} node, enforcing at the storage layer the
semantic difference between a reckless brake and a safe stop near a
pedestrian crossing. The relationship schema encodes the full causal
chain ensuring that every anomaly flagged by the analytical engine
carries a complete verifiable provenance trail back to raw sensor
readings.

% TABLE I
\begin{table}[!t]
\caption{Knowledge Graph Node Entities}
\label{tab:kg_entities}
\centering
\begin{tabular}{|p{2.2cm}|p{1.3cm}|p{4.3cm}|}
\hline
\textbf{Node Label} & \textbf{Count} & \textbf{Description} \\
\hline
Intersection & 212 &
  Campus road junction. Stores OSM id, lat/lon, betweenness and
  degree centrality (ST-GNN node features). \\
\hline
RoadSegment & 476 &
  Campus road edge. Stores length, road type and speed limit. \\
\hline
BodaBodaStop & 6 &
  Sub-label of Intersection. Known bodaboda pick-up points used
  in ContextualStop classification. \\
\hline
ZebraCrossing & 3 &
  Sub-label of Intersection. Pedestrian crossings used to
  distinguish safe stops from reckless braking. \\
\hline
Vehicle & 15 &
  Bodaboda unit with anonymised plate identifier. \\
\hline
Driver & 15 &
  Driver entity storing pseudonymised name and risk profile
  (0--1 continuous score). \\
\hline
TelematicsEvent & 826,260 &
  Enriched IoT record carrying event\_type, severity\_score,
  is\_anomaly flag, and SHA-256 provenance hash. \\
\hline
ContextualStop & -- &
  Ontologically distinct from TelematicsEvent. Represents a safe
  idle event occurring within 20~m of a BodaBodaStop or
  ZebraCrossing node. \\
\hline
ProvenanceAnchor & 2 &
  Merkle root hash for each semester batch, enabling post-ingestion
  tamper detection. \\
\hline
\end{tabular}
\end{table}

% TABLE II
\begin{table}[!t]
\caption{Knowledge Graph Relationship Types}
\label{tab:kg_relationships}
\centering
\begin{tabular}{|p{3.0cm}|p{4.8cm}|}
\hline
\textbf{Relationship} & \textbf{Description and Key Properties} \\
\hline
ADJACENT\_TO &
  Intersection $\rightarrow$ Intersection. Road connectivity edge
  carrying length\_m and road\_type. Enables 2-hop ego-graph
  traversal for causal evidence retrieval. \\
\hline
LOCATED\_AT &
  TelematicsEvent / ContextualStop $\rightarrow$ Intersection.
  HMM map-matching result with snap\_distance\_m and
  map\_match\_confidence. \\
\hline
CAUSED\_BY &
  TelematicsEvent $\rightarrow$ Driver. Links each event to the
  responsible driver, enabling driver history retrieval. \\
\hline
DRIVES &
  Driver $\rightarrow$ Vehicle. Fleet assignment relationship. \\
\hline
PRECEDES &
  TelematicsEvent $\rightarrow$ TelematicsEvent. Trajectory
  ordering edge encoding the sequential structure of each trip. \\
\hline
\end{tabular}
\end{table}

Before telemetry data is committed to the Knowledge Graph, a
privacy-preserving anonymisation pipeline is applied. Driver and
vehicle identifiers are pseudonymised using SHA-256 hashing. Raw GPS
coordinates are subjected to spatial $k$-anonymity ($k=5$, via DBSCAN
clustering) and Gaussian differential privacy noise injection
($\varepsilon=1.0$) before storage, ensuring that individual
trajectories cannot be reconstructed while aggregate safety trends
remain accurate. Access to the Knowledge Graph is governed through
Role-Based Access Control with three permission tiers---\emph{admin},
\emph{analyst}, and \emph{public}---each bound by resource-level
access control and an append-only audit log.


\section{EXPLORATORY DATA ANALYSIS}

Exploratory Data Analysis (EDA) was conducted on the 826,260-record
MakFleet simulated dataset to validate simulation parameters,
characterise class imbalance, and derive the spatial and temporal
risk schedule used to configure the ST-GNN training protocol.

\begin{figure}[H]
\centerline{\includegraphics[width=\columnwidth]{eda_temporal_hourly.png}}
\caption{Hourly trip volume (left, bar chart) and anomaly rate (right,
line plot) across the 24-hour day. Peak demand occurs at \textbf{08:00}
aligned with morning lectures. Anomaly rate is elevated in early-morning
hours (00:00--05:00) when low traffic density enables higher speeds on
empty campus roads.}
\label{fig:temporal_hourly}
\end{figure}

\textbf{Temporal Demand Analysis.}
Peak bodaboda demand occurs at 08:00, confirming the morning lecture
rush hypothesis. A secondary peak appears at 17:00 at lecture
conclusion. The anomaly rate is disproportionately elevated in
early-morning hours when low traffic density enables higher speeds on
empty roads. This bimodal pattern is captured in the ST-GNN through
timestamp features encoded in each 15-minute graph snapshot. The
semester split training protocol (L4) is motivated by this structure:
Semester~1 and Semester~2 exhibit different demand volumes due to
examination periods, representing real-world concept drift.

\begin{figure}[H]
\centerline{\includegraphics[width=\columnwidth]{eda_temporal_dow.png}}
\caption{Trip volume (left) and anomaly rate (right) by day of week.
Weekdays show higher total volume; anomaly rates remain consistent
across Monday--Friday but drop on weekends when informal campus transit
replaces schedule-driven bodaboda demand.}
\label{fig:temporal_dow}
\end{figure}

\textbf{Day-of-Week Pattern.}
Weekday trip volume substantially exceeds weekend volume. Anomaly rates
remain consistent across Monday--Friday, confirming that reckless
driving is a persistent weekday phenomenon correlated with demand
pressure. Weekend activity drops significantly, suggesting that safety
monitoring resources can be scaled down outside term-time hours.


\begin{figure}[H]
\centerline{\includegraphics[width=\columnwidth]{eda_anomaly_breakdown.png}}
\caption{Left: event type distribution across all 826,260 records.
Normal travel accounts for 91.2\%; \texttt{speeding} is the most
frequent anomaly class. Right: severity score distribution for
anomalous events showing clear separation between event types,
validating the rule-based semantic enrichment thresholds.}
\label{fig:anomaly_breakdown}
\end{figure}

\textbf{Class Imbalance and Event Type Analysis.}
Normal travel constitutes 91.2\% of the dataset while anomaly
classes---speeding, harsh braking, and rapid acceleration---account
for 8.8\%, producing an approximately 10:1 normal-to-anomaly ratio.
Speeding is the most operationally significant anomaly type. The
severity score distribution confirms that rule-based enrichment
thresholds correctly separate event types: harsh braking and speeding
cluster at severity scores above 0.7 while normal travel and idling
remain at zero. The Weighted BCE loss with $w^{+}=19.0$ addresses
the class imbalance (Section~\ref{sec:model_arch}).

\begin{figure}[H]
\centerline{\includegraphics[width=\columnwidth]{eda_speed_distribution.png}}
\caption{Box plot of vehicle speed (km/h) stratified by semantic event
type. The 30~km/h campus speed limit is shown as a dashed red line.
\texttt{speeding} events are entirely above the limit. \texttt{harsh\_braking}
spans a wide speed range confirming it is an acceleration-derived
label. \texttt{safe\_stop} and \texttt{idling} cluster near zero.
Non-overlapping distributions validate that semantic enrichment
produces meaningfully distinct speed regimes per class.}
\label{fig:speed_dist}
\end{figure}

\textbf{Speed Distribution by Semantic Event Type.}
Each ontological event type occupies a distinct speed regime,
validating the semantic enrichment pipeline. Critically, harsh braking
spans a wide speed range confirming it is an acceleration-derived
label, not a speed threshold. A sudden deceleration at 5~km/h near a
zebra crossing is a \texttt{ContextualStop}; the same signal at
40~km/h on an open road is \texttt{harsh\_braking}. Only the
spatially-aware GCN resolves this distinction.

\begin{figure}[H]
\centerline{\includegraphics[width=\columnwidth]{eda_driver_risk.png}}
\caption{Left: driver risk profiles coloured by reckless
classification (risk $>$ 0.7, red; safe, green). Right: scatter plot
of risk profile vs.\ observed anomaly rate. The positive correlation
confirms the simulation correctly links risk profile to behaviour.
Driver \textbf{Namutebi Sarah} (risk = 0.98) is the highest-risk
individual.}
\label{fig:driver_risk}
\end{figure}

\textbf{Driver Risk Profile Analysis.}
The 15 simulated drivers span the full risk spectrum from 0 to 1. The
positive correlation between assigned risk profile and observed anomaly
rate validates that the simulation correctly links driver risk to
injected anomalous behaviours. This driver-level heterogeneity is
encoded in the Knowledge Graph through the \texttt{CAUSED\_BY}
relationship, enabling the causal evidence dashboard to retrieve driver
history whenever an anomaly is flagged.

\begin{figure}[H]
\centerline{\includegraphics[width=\columnwidth]{eda_spatial_hotspots.png}}
\caption{Top-15 campus intersections by total event count (left) and
anomaly count (right). High-volume intersections cluster along the
central campus spine. The same intersections dominate both total
traffic and anomaly counts, confirming that congestion and reckless
driving are spatially co-located.}
\label{fig:spatial_hotspots}
\end{figure}

\textbf{Spatial Hotspot Analysis.}
Event density clusters along the central campus spine connecting the
main gate, Livingstone Hall, Africa Hall, and Nkurumah Hall. The same
intersections dominate anomaly counts, confirming spatial co-location
of congestion and reckless driving. Encoding each intersection as a
distinct node with betweenness and degree centrality features allows
the GCN encoder to learn that high-degree intersections carry
disproportionate risk.

\begin{figure}[H]
\centerline{\includegraphics[width=\columnwidth]{eda_semester_comparison.png}}
\caption{Event type counts (left) and proportions (right) for
Semester~1 vs.\ Semester~2. Both semesters show consistent event-type
distributions but differ in total volume. The slight shift in anomaly
proportions constitutes the concept drift the ST-GNN must generalise
across under the (L4) continual learning requirement.}
\label{fig:semester_comparison}
\end{figure}

\textbf{Semester Comparison and Concept Drift.}
Semester~1 (496,559 records) and Semester~2 (329,701 records) show
consistent event-type distributions, confirming internal coherence of
the simulation. The slight differences in anomaly proportions represent
the concept drift introduced by different academic calendar pressures.
The Temporal OOD AUC of 0.9471 (Section~\ref{sec:evaluation}) confirms
the ST-GNN successfully generalises across this drift.

\begin{figure}[H]
\centerline{\includegraphics[width=\columnwidth]{eda_sensor_validation.png}}
\caption{Left: accelerometer ($a_x$) distribution by event type,
confirming that harsh braking ($a_x < -3.5$~m/s$^2$) and rapid
acceleration ($a_x > 2.8$~m/s$^2$) thresholds separate anomaly
classes from normal driving. Right: map-matching snap distance
histogram showing 90\% of GPS points matched within 15~m. Mean
map-matching confidence = 0.906.}
\label{fig:sensor_validation}
\end{figure}

\textbf{Sensor and Map-Matching Validation.}
The accelerometer distribution confirms that the semantic enrichment
thresholds produce well-separated class distributions with minimal
overlap. The snap distance histogram validates the HMM map-matcher:
90\% of GPS points are assigned to the correct road edge within 15~m,
and mean map-matching confidence is \textbf{0.906}. Records with snap
distance above 15~m correspond to GPS readings near campus buildings
where multipath interference degrades accuracy, validating the
map-matching phase (Section~\ref{sec:methodology}) as a non-optional
preprocessing step.



\section{EXPLAINABILITY AND BUSINESS INTELLIGENCE}

The Business Intelligence (BI) layer is reconceptualised shifting from
descriptive analytics to a high-fidelity interface for Semantic
Traceability and Evidence Generation within the context of Makerere
University fleet management. This paper treats BI as a forensic window
into the Spatio-Temporal Knowledge Graph (ST-KG), answering \emph{why}
a driver was flagged rather than merely \emph{what} was recorded.

To ensure the system is actionable for non-technical stakeholders, the
BI dashboard---deployed at
\texttt{makfleet.streamlit.app}---provides three tabs: a campus
traffic heatmap, a temporal anomaly detection panel, and a Causal
Evidence tab. A Makerere University Transport Officer does not simply
view a list of risk scores; rather, the officer actively interrogates
the system's reasoning through structured causal querying.

For example, if the ST-GNN flags a high-risk anomaly for a specific
vehicle near Africa Hall, the officer selects the flagged event and
the system executes the Causal Subgraph Retrieval workflow
(Algorithm~\ref{alg:retrieval}), triggering a multi-layered evidence
generation process:
\begin{itemize}
  \item \textbf{Subgraph Extraction:} The system retrieves the
  flagged \texttt{TelematicsEvent} node, its matched
  \texttt{Intersection}, the responsible \texttt{Driver} and
  \texttt{Vehicle}, and all intersections within a 2-hop ego-graph
  via \texttt{ADJACENT\_TO} traversal.
  \item \textbf{Visual Grounding:} The BI interface renders the
  causal subgraph on a Folium campus map with colour-coded nodes
  (red = anomaly location, green = safe zones, blue = adjacent
  intersections) and generates a written explanation, for example:
  \emph{``Harsh braking detected: deceleration $-4.2$~m/s$^2$ at
  47~km/h. No zebra crossing found within 20~m in campus
  ontology---classified as reckless.''}
  \item \textbf{Driver History:} Prior anomalous events for the
  same driver are retrieved via \texttt{CAUSED\_BY} and displayed
  as an audit trail, enabling pattern-based safety interventions.
\end{itemize}

The BI interface utilises Semantic Content-Based Retrieval to bridge
the gap between the ST-GNN neural outputs and the physical reality of
the campus. Backed by the model's precision of \textbf{0.8293} on the
validation split and \textbf{0.8256} on the Semester~2 temporal
holdout, alerts suffer from minimal false positives. By allowing the
administrator to unroll the underlying causal subgraph, the system
provides a transparent logic chain from raw sensor readings back to the
high-level prediction, establishing Epistemic Trust. Consequently, any
administrative action or safety intervention is rooted in
scientifically valid, auditable evidence---transforming the BI platform
into a robust decision-support tool for institutional governance.

\section{GOVERNANCE AND PRIVACY BY DESIGN}


The high-frequency tracking of bodaboda telemetry data across the Makerere University campus inherently captures sensitive human mobility patterns. Recognising that data infrastructures are normative, the MakFleet data warehouse embeds privacy directly into its architecture. To prevent the deanonymisation of driver and passenger routines, raw GPS coordinates are subjected to spatial $k$-anonymity and differential privacy noise injection before being committed to the Knowledge Graph \cite{demontjoye2013}. This ensures that while aggregate traffic and safety trends remain highly
accurate, individual driver trajectories cannot be maliciously reconstructed. Strict Role-Based Access Control is implemented governed by university policy. While automated systems can query the graph for real-time hazard detection, full historical trajectory subgraphs are locked behind strict audit logs.
The dataset is accompanied by a formal Datasheets for Datasets artifact explicitly documenting the collection parameters, consent protocols and
potential biases in the telemetry.


\section{EVALUATION AND BENCHMARKING}
\label{sec:evaluation}

To validate the MakFleet ST-GNN, the model was evaluated under two
rigorous split protocols that enforce causal validity and prevent data
leakage. Random train/test splits are strictly forbidden because they
leak future routing behaviours into the training set.

\begin{enumerate}
  \item \textbf{Temporal Holdout Split:} The model was trained on
  Semester~1 data (January--April 2024, first 80\% of time windows)
  and evaluated on the held-out last 20\% of Semester~1 windows
  (validation) and on all Semester~2 windows (August--October 2024)
  as a temporal Out-of-Distribution (OOD) test. This directly
  measures generalisation across semester-level concept drift (L4).

  \item \textbf{Spatial Holdout Split:} Telemetry originating from
  the eastern campus zone (longitude $> 32.574°$E) was designated
  as a spatial OOD test set. The model never sees events from these
  road segments during training, directly measuring its ability to
  generalise to unseen campus topologies.
\end{enumerate}

Table~\ref{tab:temporal} reports the ST-GNN performance across all
three evaluation splits after 10 training epochs using the Adam
optimiser with Cosine Annealing scheduling.

\begin{table}[htbp]
\caption{ST-GNN Evaluation Results Across All Splits}
\label{tab:temporal}
\centering
\small
\begin{tabular}{|l|c|c|c|c|}
\hline
\textbf{Evaluation Split} & \textbf{AUC} & \textbf{F1} & \textbf{Prec.} & \textbf{Recall} \\
\hline
Validation (Sem.\,1, last 20\%) & \textbf{0.9268} & 0.9067 & 0.8293 & \textbf{1.000} \\
\hline
Temporal OOD (Semester 2) & \textbf{0.9471} & 0.9045 & 0.8256 & \textbf{1.000} \\
\hline
Spatial OOD (Eastern campus) & ---$^{\dagger}$ & \textbf{1.000} & \textbf{1.000} & \textbf{1.000} \\
\hline
\end{tabular}
\smallskip

\footnotesize{$^{\dagger}$AUC is undefined for the spatial OOD split
because all samples in the eastern campus zone are positive-class
(anomalous); F1\,=\,1.0 and Precision\,=\,1.0 confirm correct
predictions.}
\end{table}

The most significant result is that the \textbf{Temporal OOD AUC
(0.9471) exceeds the Validation AUC (0.9268)}. This counterintuitive
finding is the central academic contribution of the evaluation: it
demonstrates that the ST-GNN learned generalisable spatial risk
patterns from Semester~1 rather than memorising semester-specific
routing habits. The graph-convolutional inductive bias---propagating
risk signals through the campus road topology rather than fitting
tabular per-record features---is the direct architectural reason for
this generalisation.

\textbf{Recall = 1.000} across all three evaluation splits confirms
that the model missed zero anomalous driving events in any split. For
a campus safety application, this is the operationally critical
metric: a false negative (missed reckless driver) is far more costly
than a false positive (unnecessary alert). The Weighted BCE loss with
$w^{+}=19.0$ directly enforces this priority by asymmetrically
penalising missed anomalies.

The capability of a transport AI system to extrapolate risk policies
to unseen campus zones is critical for a scalable deployment. The
spatial OOD results confirm that the ST-GNN learns the structural
dynamics of reckless driving across varying campus topologies rather
than overfitting to the historical patterns of the central spine.

\section{DATAOPS, DEPLOYMENT \& EDGE ORCHESTRATION}


The deployment of the semantic architecture is governed by the physical constraints of the Layer 6 (L6) distributed setting, characterised by the irregular topology of the Makerere University campus and intermittent network coverage.

The system architecture is partitioned across three tiers—Edge, fog and cloud—to balance computational load and latency. These components and their respective constraints are summarised in Table~\ref{tab:physics}.

% TABLE IV
\begin{table}[!t]
\caption{Layer 6 Distributed System Physics}
\label{tab:physics}
\centering
\begin{tabular}{|l|l|p{3.8cm}|}
\hline
\textbf{Layer} & \textbf{Component} & \textbf{Constraints} \\
\hline
Edge  & IoT Devices      &
  Battery life (8--12 hrs), intermittent connectivity. \\
\hline
Fog   & Edge Gateways    &
  Limited compute (2--4 cores), batch sync every 5 min. \\
\hline
Cloud & ST-GNN Model     &
  Scalable inference, nightly model updates. \\
\hline
Cloud & Knowledge Graph  &
  ACID compliance, global replication (Neo4j). \\
\hline
Cloud & Analytics Engine &
  Serverless, event-driven processing. \\
\hline
\end{tabular}
\end{table}


To mitigate intermittent connectivity, a local SQLite cache is utilised on edge devices to store the most recent 1,000 telemetry points. A
\textit{store-and-forward} mechanism ensures that data is synchronised only when a Fog gateway is detected. Memory exhaustion at the edge is
addressed through a circular buffer and model quantisation (FP32 to INT8) which achieves a 4$\times$ reduction in the memory footprint of
the anomaly detection unit.


\section{LIMITATIONS AND ARCHITECTURAL FAILURE MODES}
\label{sec:limitations}

\textbf{Simulated Dataset.}
The MakFleet telemetry corpus is synthetically generated over the real
Makerere road network rather than collected from deployed IoT devices.
While the noise model ($\sigma = 0.00005°$, $\approx$5~m GPS drift)
and anomaly injection rates (5\% harsh braking, 3\% rapid
acceleration) are calibrated to published telematics literature, real
bodaboda behaviour introduces erratic informal-path routing and
multimodal sensor noise not fully captured in simulation. Mitigation:
deploy physical IoT telematics units on campus bodabodas and
incrementally replace simulated data with real telemetry.

\textbf{Partial ST-GNN Training.}
The ST-GNN was trained for 10 epochs on CPU. While the results are
strong (Temporal OOD AUC = 0.9471), a full 50-epoch training run on
GPU hardware would allow early stopping at the true optimum and provide
more stable OOD estimates. On an A100 GPU, the full training run is
estimated to complete in under two minutes.
Mitigation: migrate training to GPU infrastructure (e.g.\ Google Colab
Enterprise or university HPC).

\textbf{Zero ContextualStop Nodes.}
The \texttt{ContextualStop} node label --- the key ontological
distinction between safe stops and reckless braking --- was never
populated in the current simulation. Bodabodas rarely come to a full
idle within 20~m of a landmark node on the sparse simulated road
network. The label, schema, and classification logic are fully
implemented and correct; real IoT data from campus operations would
populate this label more richly.
Mitigation: increase the \texttt{SAFE\_STOP\_RADIUS\_M} threshold or
enrich the campus ontology with more granular pedestrian crossing
nodes from a field survey.

\textbf{Static Campus Road Graph.}
The campus topology is fetched once from OpenStreetMap and cached.
Road closures, new construction, or informal path changes are not
reflected until the cache is manually refreshed.
Mitigation: implement periodic OSMnx refresh triggered by a campus
facilities management API.

\textbf{GPS Accuracy Dependency.}
Signal degradation near tall campus infrastructure affects
approximately 15\% of telemetry records, requiring HMM map-matching
correction before commitment to the Knowledge Graph. In regions of
severe multipath interference, map-matching confidence drops,
increasing the risk of incorrect road-segment assignment.
Mitigation: sensor fusion with accelerometer dead-reckoning to
maintain trajectory integrity in low-GPS zones.

\section{CONCLUSION AND FUTURE WORKS}
\label{sec:conclusion}

This paper reconceptualised micro-mobility fleet management by
transitioning from antiquated relational storage to a Semantic-Aware
Spatio-Temporal Knowledge Graph. The MakFleet system was designed as
an epistemic infrastructure that enriches raw bodaboda IoT telemetry
into semantically meaningful events, stores them in an
ontology-driven Neo4j Knowledge Graph, and empowers a custom
GCN+LSTM ST-GNN to perform causally valid anomaly detection.

The central empirical finding is that the \textbf{Temporal OOD AUC of
0.9471 exceeds the Validation AUC of 0.9268}---demonstrating that the
model learned generalisable spatial risk patterns from Semester~1 that
transfer to unseen Semester~2 data, validating the (L4) continual
learning design decision. \textbf{Recall of 1.000 across all
evaluation splits} confirms that no reckless driving event was missed
in any test condition. The neuro-symbolic architecture---combining
GCN message passing with rule-based semantic event
enrichment---successfully resolves the core research challenge: a
sudden deceleration near a mapped zebra crossing is classified as a
\texttt{ContextualStop} (safe) while the identical signal on an open
road is flagged as \texttt{harsh\_braking} (anomalous), a distinction
impossible in a tabular Star Schema.

The system is publicly deployed at
\texttt{makfleet.streamlit.app}, backed by a Neo4j Aura cloud
Knowledge Graph, enabling real-time causal evidence generation for
campus safety administrators.

Future work will focus on four directions. First, deploying physical
IoT telematics units on actual Makerere campus bodabodas to replace
simulated telemetry with real-world data. Second, implementing
continual learning so the ST-GNN incrementally updates each semester
without full retraining, directly addressing (L4) concept drift.
Third, extending the campus ontology to include weather conditions,
academic calendar events, and passenger boarding zones as contextual
signals. Fourth, exploring Federated Learning so multiple Ugandan
universities can collaboratively train the model without sharing raw
driver trajectories, preserving the (L5) privacy-by-design
commitment at institutional scale.

\bibliographystyle{IEEEtran}
\clearpage
\bibliography{references}
\clearpage

\begin{IEEEbiography}[{\includegraphics[width=1in,height=1.25in,
  clip,keepaspectratio]{images/0f2a7746-253d-4089-8d77-8e614d24b522.jpeg}}]%
{Sebunya Ronaldo}
received the Bachelor of Science degree in Computer Science from Makerere University Kampala, Uganda in 2026. His undergraduate
coursework specialised in the integration of Internet of Things devices, predictive modeling and data visualisation for intelligent transport systems.

During the MakFleet research project, Mr.\ Sebunya led the development of the Internet of Things simulator and the exploratory data analysis pipeline,
producing the statistical visualisations and class-distribution analysis presented in this work.

His research interests include
predictive analytics, machine learning for anomaly detection and agentic systems.
\end{IEEEbiography}


\begin{IEEEbiography}[{\includegraphics[width=1in,height=1.25in,
  clip,keepaspectratio]{images/33333.png}}]{Ddumba Jonah}
received the Bachelor of Science degree in Computer Science from
Makerere University, Kampala, Uganda, in 2026. His academic work has focused on the design of semantic data warehousing architectures for Internet of Things-enabled transport systems and spatio-temporal knowledge graph construction.

During his undergraduate studies, Mr.\ Ddumba contributed to the MakFleet research initiative, leading the system architecture design and the integration of the Spatio-Temporal Graph Neural Network pipeline with the Makerere University campus ontology.

His research interests include data warehousing, business intelligence, Internet of Things systems and graph-based machine learning for
transport safety applications.
\end{IEEEbiography}

\begin{IEEEbiography}[{\includegraphics[width=1in,height=1.25in,
  clip,keepaspectratio]{images/IMG_5111.jpeg}}]{Ssebwana John Paul}
received the Bachelor of Science degree in Computer Science from
Makerere University, Kampala, Uganda in 2026. His undergraduate research focused on database management systems, distributed data
architectures and business intelligence pipeline design.

During his studies, Mr.\ Ssebwana contributed to the MakFleet project by designing and implementing the Knowledge Graph schema and the Role-Based Access Control governance layer for the semantic data
warehouse.

His research interests include graph database systems, data governance, privacy-preserving data engineering and scalable stream processing
architectures for real-time analytics.
\end{IEEEbiography}


\begin{IEEEbiography}[{\includegraphics[width=1in,height=1.25in,
  clip,keepaspectratio]{images/278b9526-d938-4a7a-ba58-b643e8ca6036.jpeg}}]%
{Evelyn Namuli}
received the Bachelor of Science degree in Computer Science from Makerere University, Kampala, Uganda in 2026. Her undergraduate
research concentrated on data preparation, semantic extraction and the application of natural language processing techniques to unstructured mobility data.

During the MakFleet project, Ms.\ Namuli was responsible for the LLM-assisted semantic retrieval pipeline and the differential privacy
implementation within the edge anonymisation layer.

Her research interests include semantic web technologies, naturallanguage processing, privacy-preserving data systems and the application of large language models to knowledge graph enrichment.
\end{IEEEbiography}

\EOD
\end{document}